\section{Visão integrada --- comparação final de todos os protocolos}

\preencher{Selecionar o melhor cenário de cada protocolo e explicar o critério
de seleção. Usar o arquivo large e condições comparáveis.}

\begin{table}[H]
  \centering
  \caption{Resultados consolidados das três entregas.}
  \label{tab:comparacao-final}
  \begin{tabular}{lrrrrr}
    \toprule
    Protocolo & Janela & Timeout & Throughput médio (kbps) &
      Retransmissões & Tempo total (s)\\
    \midrule
    SAW & 1  & fixo       & \preencher{kbps} & \preencher{n} & \preencher{s}\\
    GBN & 16 & fixo       & \preencher{kbps} & \preencher{n} & \preencher{s}\\
    GBN & 16 & adaptativo & \preencher{kbps} & \preencher{n} & \preencher{s}\\
    SR  & 16 & fixo       & \preencher{kbps} & \preencher{n} & \preencher{s}\\
    SR  & 16 & adaptativo & \preencher{kbps} & \preencher{n} & \preencher{s}\\
    \bottomrule
  \end{tabular}
\end{table}

\figura{comparacao_throughput.png}{0.9\linewidth}
  {Throughput médio dos protocolos e configurações comparados.}
  {fig:comparacao-throughput}

\preencher{Interpretar o gráfico de barras como uma narrativa única da evolução
do projeto, destacando ganhos, custos e o efeito do RTT adaptativo.}
